\documentclass[a4paper,twoside]{article}

\usepackage{morewrites}

\usepackage[svgnames,dvipsnames,x11names]{xcolor}
\usepackage{graphicx}
\usepackage[english]{babel}
\usepackage[colorlinks=true, linkcolor=NavyBlue, urlcolor=NavyBlue, citecolor=NavyBlue]{hyperref}
\usepackage[a4paper]{geometry}
\usepackage{ts-fonts}
\usepackage{amsmath}
\usepackage{float}
\usepackage{seqsplit}
\usepackage{ts-code}

\usepackage{lastpage}
\usepackage{refcount}
\makeatletter
\def\ps@plain{
\let\@oddhead\@empty
\let\@evenhead\@empty
\def\@oddfoot{
\hfil
\ifnum\getpagerefnumber{LastPage}>1\relax
\thepage
\fi
\hfil}
\let\@evenfoot\@oddfoot
}
\makeatother

\title{Templates}
\author{}\date{}
\begin{document}
\maketitle
Templates are predefined \LaTeX{} document structures that dictate the overall layout, style, and organization of your final document. TeXSmith includes several built-in templates for common document types, and you can also create and use custom templates.
\section{Use a template}
TeXSmith provides some templates out of the box to help you get started with common document types. You can specify a template using the \tscodeinline{-\allowbreak{}-\allowbreak{}template} (or \tscodeinline{-\allowbreak{}t}) option followed by the template name.
\begin{tscode}[lang=bash, engine=pygments]
\begin{Verbatim}[commandchars=\\\{\},breaklines, breakanywhere, commandchars=\\\{\}]
texsmith\PY{+w}{ }document.md\PY{+w}{ }\PYZhy{}\PYZhy{}template\PY{+w}{ }article
\end{Verbatim}
\end{tscode}
\section{Listing Available Templates}
To see the list of available templates, use the \tscodeinline{-\allowbreak{}-\allowbreak{}list-\allowbreak{}templates} flag:
\begin{tscode}[lang=text, stretch=0.5, engine=pygments]
\begin{Verbatim}[commandchars=\\\{\},breaklines, breakanywhere, commandchars=\\\{\}]
➜ uv run texsmith \PYZhy{}\PYZhy{}list\PYZhy{}templates
                           Available Templates
┌──────────┬─────────┬───────────────────────────────────────────────────┐
│ Name     │ Origin  │ Location                                          │
├──────────┼─────────┼───────────────────────────────────────────────────┤
│ article  │ builtin │ /home/ycr/texsmith/src/texsmith/templates/article │
│ book     │ builtin │ /home/ycr/texsmith/src/texsmith/templates/book    │
│ letter   │ builtin │ /home/ycr/texsmith/src/texsmith/templates/letter  │
│ snippet  │ builtin │ /home/ycr/texsmith/src/texsmith/templates/snippet │
└──────────┴─────────┴───────────────────────────────────────────────────┘
\end{Verbatim}
\end{tscode}
Each template comes with its own set of slots, styles, and configurations tailored for specific document types. For example, the \tscodeinline{article} template is suitable for academic papers, while the \tscodeinline{book} template is designed for longer documents with chapters and parts.
\section{Explore Template Details}
You can inspect the details of a specific template using the \tscodeinline{-\allowbreak{}-\allowbreak{}template-\allowbreak{}info} flag. This command provides information about the template's metadata, slots, attributes, and assets.
\begin{tscode}[lang=bash, engine=pygments]
\begin{Verbatim}[commandchars=\\\{\},breaklines, breakanywhere, commandchars=\\\{\}]
texsmith\PY{+w}{ }\PYZhy{}\PYZhy{}template\PY{+w}{ }article\PY{+w}{ }\PYZhy{}\PYZhy{}template\PYZhy{}info
\end{Verbatim}
\end{tscode}
It will display a summary of the template, including:
\begin{itemize}
\item List of attributes and their types (e.g., authors, columns, date, language...)
\item List of assets included in the template (e.g. style files, images...)
\item List of fragments (\tscodeinline{ts-\allowbreak{}geometry}, \tscodeinline{ts-\allowbreak{}typesetting}, ...) that the template uses.
\item List of slots available for content injection.
\end{itemize}
\section{Scaffolding Custom Templates}
One goal of TeXSmith is to make it easy to create and share custom templates. You can scaffold a new template by copying an existing one and modifying it to suit your needs. Use the \tscodeinline{-\allowbreak{}-\allowbreak{}template-\allowbreak{}scaffold} flag to copy a template to a local directory:
\begin{tscode}[lang=bash, engine=pygments]
\begin{Verbatim}[commandchars=\\\{\},breaklines, breakanywhere, commandchars=\\\{\}]
texsmith\PY{+w}{ }\PYZhy{}\PYZhy{}template\PY{+w}{ }article\PY{+w}{ }\PYZhy{}\PYZhy{}template\PYZhy{}scaffold\PY{+w}{ }my\PYZhy{}custom\PYZhy{}template/
\end{Verbatim}
\end{tscode}
Then edit the files in \tscodeinline{my-\allowbreak{}custom-\allowbreak{}template/} to customize the layout, styles, and slots as needed, and use your custom template by specifying its path:
\begin{tscode}[lang=bash, engine=pygments]
\begin{Verbatim}[commandchars=\\\{\},breaklines, breakanywhere, commandchars=\\\{\}]
texsmith\PY{+w}{ }document.md\PY{+w}{ }\PYZhy{}\PYZhy{}template\PY{+w}{ }./my\PYZhy{}custom\PYZhy{}template/
\end{Verbatim}
\end{tscode}
\end{document}
